\section{Tasks}
\label{sec:tasks}

\wcabench{} defines five tasks that span a continuous difficulty spectrum from
supervised regression to causal inference. Every task is specified by the same
five elements---formal definition, inputs and outputs, metrics, baselines and
domain challenges---and every task is exposed through a single interface so
that one evaluation code path scores all models. We use the notation
$\x$ for the model input, $\y$ for the target and $\metric$ for the metric
set.

\paragraph{The common interface.}
All tasks implement
\begin{center}\small
\code{split()} $\to$ temporal view;\quad
\code{featurize(as\_of)} $\to$ features;\quad
\code{metrics()} $\to$ metric list;\quad
\code{baseline()} $\to$ registered baselines;\quad
\code{evaluate(model)} $\to$ report.
\end{center}
The \asof{} argument is \emph{mandatory} in the feature-construction signature.
This is a design decision: leakage prevention is a type-level obligation, not
an afterthought. Baselines are organized into three tiers---trivial, domain and
methodological---and the methodological tier spans six families (statistical,
tree, deep, graph, Bayesian and causal), so that a reader can see how much is
gained over a meaningful reference point rather than over an arbitrary zero, and
can compare methodological families directly against one another.

% -----------------------------------------------------------------------------
\subsection{\task{1}: Result Prediction (Regression)}
\label{sec:task1}

\paragraph{Definition.}
Given a competitor's pre-competition history, predict the \code{best} and
\code{average} result the competitor will achieve in a specific competition,
event and round. Formally, let
$\y = f_\theta(\x)$ where
$\x = (\x^{\text{seq}}, \x^{\text{static}}, \asof)$ and
$\y = (\hat{b}, \hat{a}, \hat{\sigma}^2)$ denote predicted best, average and
predictive variance.

\paragraph{Inputs.}
Identity fields (competitor, competition, event, round type); the historical
result sequence of the last $N$ attempts; static features (age, gender,
nationality, total competition count); and the enforced decision timestamp
\asof{}.

\paragraph{Outputs.}
Continuous best and average predictions (average only where the format
supports it) and distributional parameters for interval estimation.

\paragraph{Metrics.}
Because events differ enormously in magnitude, we compute errors in the
\emph{log domain}: \maelog{} (primary) and \rmselog{} (secondary), plus
prediction-interval coverage (a calibration metric). Reporting is stratified by
competitor skill level.

\paragraph{Baselines.}
Trivial: the historical mean of the last 25 attempts. Domain/statistical:
linear extrapolation of recent results, and kernel density estimation with
bootstrap round simulation. Methodological: sequence models such as
LSTM/Transformer encoders, and gradient-boosted regression on log-transformed
targets.

\paragraph{Domain challenges.}
Performance is non-stationary: training-method improvements, injuries and
hardware changes can produce step changes in a competitor's results.
Between-event variance differs by orders of magnitude (seconds for
three-by-three, minutes for seven-by-seven, integer move counts for
fewest-moves), so event-specific normalization is essential.

% -----------------------------------------------------------------------------
\subsection{\task{2}: Placement Prediction (Ranking)}
\label{sec:task2}

\paragraph{Definition.}
Given the histories of \emph{all} competitors entered in a round, predict the
final ordering. For a round with competitor set $\mathcal{P}$, learn a scoring
function $s_\phi(\x_i)$ and output a permutation
$\hat{\pi} = \mathrm{argsort}_{i \in \mathcal{P}}(-s_\phi(\x_i))$ together with
a distribution $P(\mathrm{rank}_i = k)$ and a podium probability
$P(\mathrm{rank}_i \le 3)$.

\paragraph{Inputs.}
Competition context (competition, event, round type); history features for
every entered competitor; the enforced \asof{} timestamp.

\paragraph{Outputs.}
Predicted integer ranks, rank distributions and podium probabilities.

\paragraph{Metrics.}
Kendall's $\kendall$ rank correlation (primary), top-3 overlap with the true
podium, the Brier score of the podium probabilities, and a direct comparison
against the official \wca{} psychology sheet.

\paragraph{Baselines.}
Domain: the \wca{} \emph{psychology sheet}, which ranks competitors by their
best historical \code{ao5} and ignores variance or consistency---a strong,
practically relevant reference. Statistical: kernel-density Monte Carlo rank
simulation and the Plackett--Luce model \citep{plackett_luce}, which assigns
each competitor a latent strength and defines a distribution over orderings.
Methodological: a graph neural network over a competitor--competition
heterogeneous graph that models head-to-head interactions.

\paragraph{Domain challenges.}
The central difficulty is interaction: the same competitor may perform
differently against different fields, so the model must capture relative, not
merely absolute, strength. Round rules interact with \DNF{}: one \DNF{} is
usually discarded as the worst \code{ao5} attempt, but two \DNF{}s collapse the
round and the ranking. Round sizes vary from a handful to several hundred
competitors, which makes fixed-length ranking architectures awkward.

% -----------------------------------------------------------------------------
\subsection{\task{3}: DNF Prediction (Imbalanced Classification)}
\label{sec:task3}

\paragraph{Definition.}
Given the attempts a competitor has already completed in a round and the
number remaining, predict the probability that a subsequent attempt is a
\DNF{}. Formally, with partial sequence
$\x = (x_1, \dots, x_{k})$ and round context $c$ (format, attempts left),
estimate $P(\text{DNF} \mid \x, c)$ and the whole-round probability
$P(\text{round DNF})$.

\paragraph{Inputs.}
Competitor identity; the completed attempt values (including \DNF{} markers);
round format and remaining attempts; the competitor's historical \DNF{} rate;
the enforced \asof{} timestamp.

\paragraph{Outputs.}
A per-attempt \DNF{} probability and a whole-round \DNF{} probability.

\paragraph{Metrics.}
\aucroc{} and \aucpr{} (with the latter emphasized under imbalance), F1, the
Matthews correlation coefficient (MCC), and a reliability diagram for
calibration. Because \aucroc{} flatters models when positives are rare, we
require \aucpr{} and MCC alongside a random/global-rate reference. We also
report a \emph{hard subset} on which the competitor's historical \DNF{} rate
lies in $[0.1, 0.3]$, to preserve discriminative power against saturation.

\paragraph{Baselines.}
Trivial/domain: the competitor's historical \DNF{} rate. Statistical:
logistic regression on completed-attempt mean and variance, and a Bayesian
hierarchical model with a Beta prior for shrinkage. Methodological: random
forest / gradient boosting over competitor and round-context features.

\paragraph{Domain challenges.}
\DNF{} is rare and unevenly distributed: blindfolded and multi-blind events
fail far more often, while three-by-three is stable, and fewest-moves has its
own timeout/violation failure mode. \DNF{} occurrence is also
sequence-dependent: after a failure a competitor may become more conservative
(reducing risk) or may lose composure (increasing it). Finally, risk is
rule-dependent---after one \DNF{} in \code{ao5} the tolerance for another
failure is zero.

% -----------------------------------------------------------------------------
\subsection{\task{4}: Human-Limit Estimation (Extreme-Value Extrapolation)}
\label{sec:task4}

\paragraph{Definition.}
From each event's historical world-record series
$\{(t_i, y_i)\}$ estimate the theoretical performance ceiling $L$ and the
year at which it is expected to be approached. Employ the saturating model
\begin{equation}
  y(t) = L + (y_0 - L)\,\exp\!\big(-\lambda\, g(t)\big) + \varepsilon(t),
  \label{eq:limit}
\end{equation}
where $L$ is the limit, $g(t)$ is a technological-progress function and
$\varepsilon$ is noise, and infer the posterior $p(L, \lambda, \text{params}
\mid \text{data})$.

\paragraph{Inputs.}
The world-record series per event (date, value), plus event-level features
(mean, variance, history length).

\paragraph{Outputs.}
A point estimate and posterior for $L$, the predictive interval, and a
convergence-year estimate with an interval.

\paragraph{Metrics.}
There is no observable ground truth for a limit, so evaluation relies on
(a) stability of the estimate under leave-one-out refitting, (b) empirical
coverage of the predictive intervals, and (c) consistency with established
domain knowledge. As an anchor, the speedcubing community estimates the
three-by-three limit near $2.37$ seconds with convergence around 2036.

\paragraph{Baselines.}
Statistical: the exponential-decay model of Equation~\ref{eq:limit}.
Methodological: Gaussian-process regression combined with extreme-value theory
\citep{evt_coles}; a Bayesian hierarchical extreme-value model
that treats per-event limits as draws from a hyper-distribution
\citep{bayesian_hierarchical}; and change-point detection that identifies
structural breaks in the rate of improvement.

\paragraph{Domain challenges.}
The dominant risk is extrapolation: historical data cannot anticipate a future
jump caused by new hardware or training regimes. Data density varies wildly;
three-by-three has 23 years of dense data, whereas some events provide fewer
than fifty record points and demand strong priors or hierarchical borrowing.
Finally, the target is unobservable, so validation must be indirect.

% -----------------------------------------------------------------------------
\subsection{\task{5}: Skill Transfer (Causal Inference)}
\label{sec:task5}

\paragraph{Definition.}
Quantify how an improvement in one event causally affects performance in
another. For each ordered event pair $(A, B)$ estimate the causal effect
$\theta_{A \to B}$ and assemble a $17 \times 17$ transfer matrix with
confidence intervals.

\paragraph{Inputs.}
Competitors' multi-event records (timestamp, event, result) and the time at
which each competitor first participated in each event.

\paragraph{Outputs.}
A transfer matrix $\Theta = [\theta_{A \to B}]$ and interval estimates, with an
explicit identifiability flag for pairs that cannot be estimated reliably.

\paragraph{Metrics.}
Point-estimate accuracy against held-out data, robustness of the estimate
across subsamples, and agreement with domain-expert judgement. Because a
ground-truth effect is unavailable, we additionally report the
\emph{number of identifiable event pairs}.

\paragraph{Baselines.}
Correlational: Pearson/Spearman correlation between event performances---the
naive approach that ignores confounding. Causal: difference-in-differences
(DID) \citep{did} treating first participation in event $A$ as treatment;
instrumental variables (IV) \citep{iv} using the first domestic hosting of an
event as an instrument; and a causal forest \citep{causal_forest} for
heterogeneous treatment effects.

\paragraph{Domain challenges.}
Confounding is severe: talent and training effort raise performance in all
events simultaneously, producing spurious positive transfer. Participation
order is non-random---competitors who are good at event $A$ may take up event
$B$ earlier---which induces self-selection bias. Finally, many event pairs have
very few common competitors, so effects are poorly identified and the benchmark
requires an explicit ``not identifiable'' label rather than a forced estimate.
